\documentclass[../sablona-main.tex]{subfiles}
\begin{document}

\section{Řešení modelu simplexovou metodou}
Pro vyřešení modelu simplexovým algoritmem musíme úlohu převést z kanonického tvaru na tvar standardní. K původním čtyřem strukturním omezením (nohy, dřevěné desky, skleněné desky a montážní čas) zavedeme nezáporné doplňkové proměnné $s_1, s_2, s_3, s_4 \ge 0$. Tyto proměnné vyrovnávají nerovnosti na rovnice a reprezentují nevyužitou kapacitu jednotlivých zdrojů.

\subsection{Standardní tvar modelu}
Cílovou funkci maximalizace zisku i omezující podmínky zapíšeme v souladu s obecnou notací modelu:

\begin{equation}
	\max Z = 200x_1 + 350x_2 + 0 \cdot s_1 + 0 \cdot s_2 + 0 \cdot s_3 + 0 \cdot s_4
\end{equation}

za podmínek:
\begin{align}
	5x_1 + 5x_2 + s_1 &= 300 \quad \text{(nohy)} \\
	x_1 + s_2 &= 50 \quad \text{(dřevo)} \\
	x_2 + s_3 &= 35 \quad \text{(sklo)} \\
	0.6x_1 + 1.5x_2 + s_4 &= 63 \quad \text{(montáž)}
\end{align}
přičemž $x_1, x_2, s_1, s_2, s_3, s_4 \ge 0$.

\subsection{Průběh výpočtu a interpretace}
Počáteční bazické řešení odpovídá stavu, kdy se nic nevyrábí ($x_1=0, x_2=0$) a veškeré kapacity jsou volné. Algoritmus v každém kroku hledá hranu konvexního polyedru, podél které účelová funkce nejrychleji roste.

\begin{itemize}
	\item Vstupující proměnná: Jako první vstupuje do báze $x_2$ (luxusní stoly), protože má nejvyšší jednotkový zisk \$350.
	\item Klíčový prvek: Omezením pro $x_2$ je buď počet skleněných desek (max 35 ks), nebo montážní čas (max 42 ks). Rozhoduje nejpřísnější omezení.
	\item Iterace: Po provedení pivotových operací docházíme k optimálnímu řešení v bodě $[30, 30]$.
\end{itemize}

\subsection{Optimální výsledky}
Simplexová metoda potvrdila výsledky získané grafickým řešením a softwarovými nástroji Excel Solver a GAMS:
\begin{itemize}
	\item \textbf{Počet základních stolů ($x_1$):} 30 ks 
	\item \textbf{Počet luxusních stolů ($x_2$):} 30 ks
	\item \textbf{Maximální celkový zisk ($Z$):} \$16\,500 
\end{itemize}

Z hlediska analýzy zdrojů zůstávají doplňkové proměnné $s_2 = 20$ (nevyužité dřevěné desky) a $s_3 = 5$ (nevyužité skleněné desky) nenulové. Zdroje nohy a montážní čas jsou vyčerpány beze zbytku ($s_1=0, s_4=0$), což z nich činí kritická omezení výroby.
\end{document}
